\scanpage{4}
donc sont isomorphes. Il reste à voir que pour tt $\Sigma$,
\[
\operatorname{Aut}(\Sigma)\xrightarrow{\unsure{\prime|\sim}}\operatorname{Aut}(\Phi(\Sigma))
\]
est un iso de groupes. Or on a un hom. canonique
\[
W\xrightarrow{\alpha}\operatorname{Aut}(\Phi(\Sigma))
\]
par
\[
\alpha(w)=(\operatorname{int}(w),w_R,\operatorname{id}_I)
\]
[en effet,
\[
\begin{aligned}
w_R(\deleted{\unsure{w}}\,\underset{\substack{\uparrow\\W}}{w'}\underset{\substack{\uparrow\\R}}{r})&=\operatorname{int}(w)(w')\,w_R\,r\\
\Vert\qquad &\qquad\Vert\\
w(w'r)&\qquad(ww'w^{-1})wr
\end{aligned}
\]
et
\[
\left.\begin{aligned}
\varphi(w_R(r)\deleted{\unsure{r}},i)&=\operatorname{int}(w)\varphi(r,i)\\
\Vert\qquad &\\[-2pt]
\varphi(wr,i)&
\end{aligned}\right\}\ \text{par a)}
\]
Ceci dit, $\alpha$ est \unsure{évidemt|évidemment} injectif, son image est formé des autom. \deleted{\unsure{$(u_W,u_R$|$(u_W,u_R)$}} de $\Phi(\Sigma)=(W,R,I,\varphi)$ induisant l'identité sur $I$, i.e. de la forme $(u_W,u_R,\operatorname{id}_I)$, où $u_R$ est un autom affine de $R$, associé à l'autom $u_W$ du grpe $W$ (la \textquotedblleft partie homogène\textquotedblright\ de $u_R$) et satisfaisant
\[
\begin{aligned}
\unsure{(*)|(\alpha)|(\ast\ast)}\quad \deleted{\unsure{*}}\qquad
\varphi(u_Rr,i)&=u_W\,\varphi(r,i)\quad\text{i.e.}\\
\varphi(r,i)&=u_W^{-1}\,\varphi(u_Rr,i)
\end{aligned}
\]
Choisissons en effet $r_0\in R$ [\deleted{\unsure{on aura tout}} \inserted{permettant d'identifier} $R$ à $W$] donc on aura
\[
u_R(wr_0)=u_W(w)ar_0\quad\text{où }\deleted{\unsure{a}}u_R\unsure{(r_0)|wr_0)}=ar_0
\]
et la condition \unsure{$(*)$|$(\alpha)$|$(\ast\ast)$} \deleted{\unsure{a}} s'écrit (pour $r=wr_0\sim w$
% archive sheet number: 3

